\section{Why Unplanned Intentions Often Fail}

\begin{proposition}[Unplanned intention preserves current advantage]
Assume two future options compete at a decision window: remaining in the
present state $s_0$ or attempting transition to $s_1$. If no plan lowers the
barrier terms for $s_1$, then the model provides no new structural reason for
$s_1$ to become easier relative to the status quo at that window.
\end{proposition}

\begin{proof}
By assumption, no plan lowers ambiguity, raises preparation, or otherwise
reduces the transition burden for $s_1$. Therefore the model's barrier-relevant
variables are unchanged in favour of $s_1$. Since the current state's
structural advantage has not been reduced by planning, the model introduces no
new formal reason for the transition to become easier relative to the status
quo.
\end{proof}

\begin{corollary}[Unplanned preference alone does not reverse inertia]
Within this framework, a merely verbal preference for $s_1$ does not by itself
imply that $s_1$ will defeat the current state's structural advantage at the
next decision window.
\end{corollary}

\begin{proof}
Apply the previous proposition to the special case where the only new element
is verbal endorsement, and no barrier-relevant planning change has occurred.
\end{proof}

\begin{example}[Exercise after work]
The sentence ``I should work out after work'' is not yet a usable plan. It
leaves open the location, bag, sequence, exercise selection, session length,
and exit condition. In the chapter's notation, these unresolved items keep
execution complexity high and leave target ambiguity largely unchanged. A real
plan resolves some of those variables upstream.
\end{example}
